\chapter{Tinjauan Pustaka dan Dasar Teori}

\section{Tinjauan Pustaka}

Pemantauan kerumunan berbasis deteksi objek menggunakan model YOLO telah menjadi topik 
penelitian yang aktif dalam beberapa tahun terakhir. Sebagian penelitian berfokus pada 
penghitungan jumlah orang secara waktu nyata dari rekaman kamera. Gündüz dan Işık 
\cite{gunduz2023new} mengusulkan metode deteksi kerumunan waktu nyata yang menentukan 
batas suatu wilayah pada video, menghitung jumlah orang di dalamnya, dan membandingkan 
performa beberapa varian model YOLO untuk keperluan tersebut. Pendekatan dengan tujuan 
serupa diterapkan Ginting dkk. \cite{ginting2021crowd} yang menggabungkan YOLOv3-Tiny dan 
algoritma Viola-Jones untuk mendeteksi kerumunan di pusat perbelanjaan serta memicu 
peringatan ketika jumlah orang melampaui ambang tertentu. Pada penerapan di perangkat 
dengan sumber daya terbatas, Hansika dkk. \cite{hansika2024privacy} membangun kerangka 
penghitungan kerumunan ringan menggunakan YOLOv8n dan algoritma pelacakan Norfair yang 
dijalankan pada perangkat tepi Raspberry Pi 5, dengan hasil ditampilkan melalui 
\textit{dashboard}.

Sebagian penelitian lain menekankan keterpaduan deteksi dengan antarmuka pemantauan dan 
sistem peringatan. Apeksha dkk. \cite{apeksha2025transit} mengembangkan platform 
AI-Transit Alerts yang menghitung jumlah penumpang transportasi umum secara waktu nyata 
menggunakan YOLOv8, menampilkannya pada \textit{dashboard} interaktif, dan mengirim 
peringatan otomatis ketika jumlah penumpang melampaui batas aman. Wong dan Law 
\cite{wong2024fusion} menyoroti keterbatasan pemantauan CCTV manual yang melelahkan dan 
rawan kesalahan pada ruang publik padat, lalu mengusulkan kerangka yang menggabungkan 
rekaman CCTV dengan informasi denah ruang untuk memetakan lintasan individu dan memantau 
kepadatan secara waktu nyata.

Di luar deteksi orang, model YOLO juga digunakan untuk deteksi kendaraan dalam pemantauan 
lalu lintas. Zhang dkk. \cite{zhang2022realtime} mengembangkan model YOLOv5 yang 
ditingkatkan untuk deteksi kendaraan waktu nyata dengan fokus menurunkan tingkat deteksi 
keliru (\textit{false detection}). Ringkasan perbandingan penelitian terdahulu disajikan 
pada Tabel \ref{tab:tinjauan}.

\begin{table}[h]
    \centering
    \caption{Perbandingan penelitian terdahulu.}
    \begin{tabular}{|p{2.4cm}|p{2.6cm}|p{3.0cm}|p{3.6cm}|}
        \hline
        \textbf{Penelitian} & \textbf{Model} & \textbf{Domain Aplikasi} & \textbf{Catatan} \\
        \hline \hline
        Gündüz \& Işık \cite{gunduz2023new} & Beberapa varian YOLO & Penghitungan kerumunan dalam ruang & Membandingkan varian YOLO dan menentukan wilayah hitung \\
        \hline
        Ginting dkk. \cite{ginting2021crowd} & YOLOv3-Tiny, Viola-Jones & Deteksi kerumunan di pusat perbelanjaan & Memicu peringatan saat ambang terlampaui \\
        \hline
        Hansika dkk. \cite{hansika2024privacy} & YOLOv8n, Norfair & Penghitungan kerumunan pada perangkat tepi & Ringan, hemat sumber daya, dengan \textit{dashboard} \\
        \hline
        Apeksha dkk. \cite{apeksha2025transit} & YOLOv8 & Penghitungan penumpang transportasi umum & \textit{Dashboard} interaktif dan peringatan otomatis \\
        \hline
        Wong \& Law \cite{wong2024fusion} & Model \textit{deep learning} & Pemantauan kepadatan ruang publik & Menggabungkan CCTV dengan denah ruang \\
        \hline
        Zhang dkk. \cite{zhang2022realtime} & YOLOv5 (ditingkatkan) & Deteksi kendaraan lalu lintas & Menurunkan tingkat deteksi keliru \\
        \hline
        Penelitian ini & YOLOv8 \& YOLO11 (n, s, m, l) & Deteksi sembilan kelas objek di Malioboro & Multi-kelas mencakup andong, becak, dan bajaj. Membandingkan dua generasi dan diimplementasikan sebagai sistem berbasis web \\
        \hline
    \end{tabular}
    \label{tab:tinjauan}
\end{table}

Penelitian-penelitian tersebut menunjukkan model YOLO efektif untuk pemantauan kerumunan 
maupun kendaraan secara waktu nyata. Meskipun demikian, sebagian besar berfokus pada satu 
jenis objek saja, yaitu orang atau kendaraan, dan mengandalkan kelas objek yang tersedia 
pada dataset umum. Belum ada penelitian yang menangani kombinasi objek khas kawasan 
wisata Indonesia, termasuk andong, becak, dan bajaj, dalam satu sistem. Perbandingan 
antar varian model pada sebagian besar penelitian juga terbatas pada satu generasi YOLO. 
Penelitian ini menempatkan diri pada celah tersebut dengan membandingkan varian YOLOv8 
dan YOLO11 untuk mendeteksi sembilan kelas objek khas kawasan Malioboro secara bersamaan, 
kemudian mengimplementasikan model terbaik ke dalam sistem pemantauan berbasis web.

\section{Dasar Teori}

Subbab ini menyusun landasan teori mengikuti alur kerja sistem yang dibangun, mulai dari 
teori deteksi objek hingga penerjemahan hasil deteksi menjadi penilaian kepadatan. Teori 
deteksi objek dan arsitektur YOLO menjelaskan cara objek dikenali dari citra kamera 
pengawas, sedangkan metrik evaluasi menetapkan ukuran kualitas deteksi yang dihasilkan. 
Ketidakseimbangan kelas dan anotasi semi-otomatis membahas dua persoalan data yang 
menentukan keberhasilan pelatihan model pada konteks kawasan Malioboro. Bagian terakhir, 
yaitu tingkat pelayanan pejalan kaki, menghubungkan keluaran model dengan tujuan akhir 
sistem berupa estimasi kepadatan dan penetapan ambang notifikasi.

\subsection{Deteksi Objek}

Deteksi objek adalah tugas \textit{computer vision} yang bertujuan mengidentifikasi 
lokasi dan kategori setiap objek yang ada dalam suatu gambar secara bersamaan 
\cite{oksuz2021imbalance}. Keluaran dari sistem deteksi objek berupa kotak pembatas 
(\textit{bounding box}) yang melingkupi setiap objek beserta label kelas dan nilai 
keyakinan (\textit{confidence score}) yang menyertainya.

Pendekatan deteksi objek berbasis \textit{deep learning} secara umum dibagi menjadi dua 
kategori, yaitu \textit{two-stage detector} dan \textit{one-stage detector}. 
\textit{Two-stage detector} bekerja dalam dua tahap terpisah: tahap pertama menghasilkan 
kandidat wilayah yang mungkin mengandung objek, dan tahap kedua mengklasifikasikan serta 
menyempurnakan posisi kotak pembatas pada setiap kandidat wilayah tersebut. Faster R-CNN 
\cite{ren2015faster} adalah salah satu arsitektur \textit{two-stage} yang berpengaruh, 
yang memperkenalkan \textit{Region Proposal Network} (RPN) untuk menghasilkan proposal 
wilayah secara efisien dengan berbagi fitur konvolusional antara tahap proposal dan tahap 
deteksi. Pendekatan ini menghasilkan akurasi yang tinggi, namun kecepatan inferensinya 
terbatas karena sifat dua tahapnya, yaitu sekitar 5 \textit{frame} per detik pada GPU.

\textit{One-stage detector} menggabungkan kedua tahap tersebut menjadi satu jaringan 
tunggal yang langsung memprediksi kotak pembatas dan label kelas dari seluruh gambar 
dalam satu kali evaluasi \cite{redmon2016yolo}. Pendekatan ini jauh lebih cepat 
dibandingkan \textit{two-stage detector} dan lebih sesuai untuk kebutuhan deteksi waktu 
nyata, meskipun pada generasi awal memiliki akurasi yang sedikit lebih rendah. Penelitian 
ini menggunakan pendekatan \textit{one-stage} karena kebutuhan pemantauan kepadatan 
secara waktu nyata dari aliran video kamera pengawas.

Komponen penting dalam arsitektur detektor modern adalah \textit{Feature Pyramid Network} 
(FPN) \cite{lin2017fpn}. FPN membangun representasi piramida fitur multi-skala dari 
jaringan konvolusional dengan biaya komputasi yang minimal, menggunakan arsitektur 
\textit{top-down} dengan koneksi lateral untuk menghasilkan peta fitur berkualitas tinggi 
di semua skala. Kemampuan mendeteksi objek dalam berbagai ukuran ini menjadikan FPN 
komponen penting dalam detektor YOLO generasi terbaru, termasuk YOLOv8 dan YOLO11 yang 
digunakan dalam penelitian ini.

\subsection{Arsitektur YOLO}

YOLO (\textit{You Only Look Once}) diperkenalkan oleh Redmon dkk. \cite{redmon2016yolo} 
sebagai pendekatan deteksi objek \textit{one-stage} yang memformulasikan deteksi sebagai 
masalah regresi tunggal dari piksel gambar ke koordinat kotak pembatas dan probabilitas 
kelas. Model dasar YOLO mampu memproses gambar pada kecepatan 45 \textit{frame} per detik, 
sedangkan versi yang lebih ringan mencapai 155 \textit{frame} per detik. Sejak 
diperkenalkan, arsitektur YOLO telah berkembang melalui banyak iterasi yang membawa 
peningkatan pada akurasi, efisiensi komputasi, dan kemampuan multitugas 
\cite{terven2023comprehensive}.

YOLOv8 dirilis oleh Ultralytics pada Januari 2023 dan tersedia dalam lima ukuran, yaitu 
nano (n), small (s), medium (m), large (l), dan extra-large (x) 
\cite{terven2023comprehensive}. Arsitektur YOLOv8 memperkenalkan modul C2f 
(\textit{cross-stage partial bottleneck with two convolutions}) pada \textit{backbone} 
sebagai pengganti modul C3 pada versi sebelumnya. C2f menggabungkan fitur tingkat tinggi 
dengan informasi kontekstual untuk meningkatkan akurasi deteksi. YOLOv8 juga menggunakan 
model \textit{anchor-free} dengan \textit{decoupled head} yang memproses tugas klasifikasi 
dan regresi kotak pembatas secara terpisah, sehingga mengurangi konflik antar-tugas dan 
meningkatkan konvergensi pelatihan.

YOLO11 dirilis oleh Ultralytics pada September 2024 sebagai penerus YOLOv8 
\cite{khanam2024yolov11}. Khanam dan Hussain menganalisis inovasi arsitektur YOLO11, yang 
mencakup blok C3k2 (\textit{Cross Stage Partial with kernel size 2}), komponen SPPF 
(\textit{Spatial Pyramid Pooling Fast}), dan modul C2PSA 
(\textit{Convolutional block with Parallel Spatial Attention}). Blok C3k2 meningkatkan 
ekstraksi fitur dengan kernel berukuran lebih kecil yang lebih efisien secara komputasi, 
sedangkan C2PSA menambahkan mekanisme perhatian spasial paralel untuk memperkuat kemampuan 
model menangkap konteks global dalam gambar. YOLO11 tersedia dalam skala nano hingga 
extra-large dan mendukung berbagai tugas \textit{computer vision} termasuk deteksi objek, 
segmentasi, dan estimasi pose.

Kedua model tersebut, yaitu YOLOv8 dan YOLO11, dipilih dalam penelitian ini karena 
keduanya merupakan model resmi Ultralytics yang berbagi kerangka pelatihan, format data, 
dan antarmuka pemrograman yang identik. Keseragaman ini memungkinkan perbandingan yang 
terkontrol, di mana perbedaan performa dapat diatribusikan pada perbedaan arsitektur dan 
bukan pada perbedaan implementasi. Selain itu, kedua model memiliki skalabilitas yang 
baik untuk perangkat dengan sumber daya terbatas, yang relevan untuk kebutuhan inferensi 
waktu nyata pada sistem pemantauan \cite{terven2023comprehensive}.

Model transformer seperti RT-DETR \cite{zhao2024rtdetr} merupakan alternatif yang 
kompetitif dari segi akurasi, namun YOLOv8 dan YOLO11 dipilih karena efisiensi parameter, 
ekosistem yang matang, serta kemampuan deployment yang lebih luas pada berbagai jenis 
perangkat mulai dari server hingga perangkat tepi.

\subsection{Metrik Evaluasi Deteksi Objek}

Evaluasi performa model deteksi objek menggunakan sejumlah metrik standar yang saling 
melengkapi \cite{padilla2020survey}. Metrik-metrik ini digunakan secara konsisten pada 
seluruh eksperimen dalam penelitian ini.

\textit{Intersection over Union} (IoU) mengukur seberapa besar tumpang tindih antara 
kotak pembatas hasil prediksi dan kotak pembatas acuan (\textit{ground truth}). IoU 
dihitung sebagai rasio luas irisan terhadap luas gabungan kedua kotak. Sebuah prediksi 
dinyatakan benar (\textit{true positive}) jika nilai IoU-nya melampaui ambang tertentu.

\textit{Precision} mengukur proporsi prediksi positif yang benar dari seluruh prediksi 
yang dihasilkan model, sedangkan \textit{recall} mengukur proporsi objek yang berhasil 
dideteksi dari seluruh objek yang sebenarnya ada. Kedua metrik ini dihitung per kelas.

\textit{Average Precision} (AP) merangkum kinerja detektor dengan menghitung luas area 
di bawah kurva \textit{precision-recall} \cite{padilla2020survey}. AP dihitung secara 
terpisah untuk setiap kelas objek. \textit{mean Average Precision} (mAP) adalah rata-rata 
AP di seluruh kelas yang dievaluasi.

Dataset COCO \cite{lin2014coco} memperkenalkan dua varian mAP yang digunakan secara luas 
sebagai standar evaluasi. mAP50 dihitung pada ambang IoU 0,50, sementara mAP50-95 
merupakan rata-rata mAP pada ambang IoU dari 0,50 hingga 0,95 dengan langkah 0,05. 
mAP50-95 memberikan evaluasi yang lebih ketat karena memperhitungkan ketepatan lokalisasi 
kotak pembatas pada berbagai tingkat kesesuaian. Dataset COCO sendiri memuat anotasi 91 
jenis objek pada 328 ribu gambar dengan total 2,5 juta \textit{instance} berlabel, dan 
telah menjadi tolok ukur standar untuk evaluasi detektor objek.

\textit{F1-score} adalah rata-rata harmonik dari \textit{precision} dan \textit{recall}, 
memberikan metrik tunggal yang menyeimbangkan kedua aspek tersebut. F1 sangat berguna 
untuk mengevaluasi kinerja per kelas pada kondisi ketidakseimbangan kelas, karena 
memberikan penalti yang setara terhadap \textit{precision} maupun \textit{recall} yang 
rendah.

Dalam penelitian ini, metrik utama yang digunakan adalah mAP50, mAP50-95, 
F1-\textit{score} per kelas, dan \textit{confusion matrix} untuk menganalisis pola 
kesalahan deteksi.

\subsection{Ketidakseimbangan Kelas pada Deteksi Objek}

Ketidakseimbangan kelas adalah kondisi ketika jumlah \textit{instance} antar kelas dalam 
dataset berbeda jauh. Oksuz dkk. \cite{oksuz2021imbalance} mengidentifikasi 
ketidakseimbangan kelas sebagai salah satu masalah utama dalam deteksi objek, yang dapat 
menyebabkan model lebih condong mempelajari kelas dengan jumlah \textit{instance} besar 
dan mengabaikan kelas dengan jumlah \textit{instance} kecil. Dampaknya terlihat pada 
\textit{Average Precision} yang rendah untuk kelas minoritas meskipun mAP keseluruhan 
tampak cukup baik.

Dataset yang dibangun dalam penelitian ini menunjukkan ketidakseimbangan kelas yang 
ekstrem. Kelas orang memiliki 10.694 \textit{instance} pada data latih, sedangkan kelas 
truk hanya memiliki 45 \textit{instance}, dengan rasio sebesar 237:1. Kelas bajaj dan bus 
masing-masing hanya memiliki 96 dan 92 \textit{instance}. Kondisi ini merupakan 
konsekuensi langsung dari komposisi objek yang ada di kawasan Malioboro, bukan kesalahan 
pengumpulan data.

Ketidakseimbangan ini diperkirakan berpengaruh pada performa per kelas model yang dilatih, 
khususnya pada kelas dengan \textit{instance} sedikit seperti truk, bajaj, dan sepeda. 
Analisis pengaruh ketidakseimbangan ini terhadap F1-\textit{score} per kelas disajikan 
pada BAB IV.

\subsection{Anotasi Semi-Otomatis}

Anotasi dataset dalam jumlah besar secara manual membutuhkan waktu dan sumber daya yang 
besar. Pendekatan anotasi semi-otomatis menggabungkan prediksi model \textit{pretrained} 
dengan verifikasi dan koreksi manual untuk mempercepat proses anotasi tanpa mengorbankan 
kualitas secara keseluruhan \cite{tenckhoff2025feedback}. Tenckhoff dkk. menunjukkan 
bahwa platform anotasi berbasis umpan balik model, di mana model \textit{pretrained} 
menghasilkan anotasi awal yang kemudian diverifikasi dan diperbaiki oleh anotator manusia, 
menghasilkan penghematan waktu dan jumlah interaksi anotasi yang terukur dibandingkan 
anotasi penuh secara manual.

Dalam penelitian ini, pendekatan ini diterapkan pada 1.252 gambar yang dianotasi 
menggunakan model \textit{pretrained} untuk enam kelas yang tersedia pada dataset umum, 
diikuti dengan verifikasi dan koreksi manual. Kelas yang tidak tersedia pada model 
\textit{pretrained}, yaitu andong dan becak, dianotasi secara penuh manual. Himpunan 
validasi dan pengujian sebanyak 329 gambar dianotasi secara penuh manual untuk memastikan 
integritas evaluasi. Rincian pipeline anotasi diuraikan pada BAB III.

\subsection{Tingkat Pelayanan Pejalan Kaki}

Tingkat pelayanan (\textit{level of service}, LOS) pejalan kaki adalah konsep kuantitatif 
yang mengklasifikasikan kondisi pergerakan pejalan kaki berdasarkan kepadatan dan 
kenyamanan. Fruin \cite{fruin1971designing} mendefinisikan enam tingkat pelayanan 
(A hingga F) berdasarkan studi fotografi \textit{time-lapse} terhadap pergerakan pejalan 
kaki, dengan parameter utama berupa kepadatan orang per meter persegi dan ruang yang 
tersedia per pejalan kaki. Tingkat A menggambarkan kondisi pergerakan bebas tanpa konflik 
antar-pejalan kaki, sedangkan tingkat E dan F menggambarkan kondisi padat di mana kontak 
fisik tidak terhindarkan dan pergerakan mandiri menjadi terbatas.

Di Indonesia, konsep LOS ini diadopsi dalam Peraturan Menteri PUPR Nomor 03/PRT/M/2014 
tentang Pedoman Perencanaan, Penyediaan, dan Pemanfaatan Prasarana dan Sarana Jaringan 
Pejalan Kaki di Kawasan Perkotaan \cite{pupr2014pedoman}, yang menjadi acuan standar 
nasional untuk perencanaan jalur pedestrian. Penerapan standar nasional ini memberikan 
dasar ilmiah yang terukur untuk penetapan ambang notifikasi kepadatan pada sistem 
pemantauan yang dibangun dalam penelitian ini, khususnya untuk kelas orang pada kamera 
yang dikalibrasi.

Untuk menghitung kepadatan dari hasil deteksi, diperlukan pemetaan antara koordinat 
piksel dan koordinat nyata di bidang tanah. Pemetaan ini dilakukan melalui estimasi 
homografi planar, yaitu transformasi perspektif yang mengubah titik-titik pada bidang 
gambar menjadi koordinat metrik di bidang tanah menggunakan korespondensi titik acuan 
yang diketahui dimensinya \cite{fruin1971designing}. Dengan homografi yang telah 
dikalibrasi, jumlah objek yang terdeteksi dalam suatu zona dapat dibagi dengan luas zona 
tersebut dalam meter persegi untuk memperoleh estimasi kepadatan, yang kemudian dipetakan 
ke tingkat pelayanan sesuai klasifikasi Fruin dan Permen PUPR 2014.

\section{Analisis Perbandingan Metode}

\label{sec:analisis_perbandingan}

Pemilihan metode deteksi objek untuk sistem pemantauan di kawasan Malioboro tidak dapat 
dilepaskan dari kebutuhan spesifik permasalahan yang dihadapi. Sistem harus melakukan 
inferensi secara \textit{real-time} pada perangkat dengan sumber daya terbatas, mampu 
mendeteksi pejalan kaki sekaligus kelas kendaraan lokal yang tidak tercakup dalam model 
umum, dapat dilatih dari \textit{dataset} kustom berukuran terbatas, serta layak 
diterapkan pada lingkungan operasional. Empat kebutuhan inilah yang menjadi tolok ukur 
dalam membandingkan metode-metode yang tersedia. Subbab ini meninjau kelebihan dan 
kekurangan tiap pendekatan terhadap tolok ukur tersebut, lalu menjustifikasi metode yang 
dipilih.

\subsection{Perbandingan Pendekatan Deteksi Objek}
\label{subsec:perbandingan_pendekatan}

Pendekatan \textit{two-stage} yang diwakili oleh Faster R-CNN memisahkan deteksi menjadi 
dua tahap, yaitu pembangkitan usulan wilayah (\textit{region proposal}) dan klasifikasi 
pada tiap usulan tersebut \cite{ren2015faster}. Pemisahan ini menghasilkan akurasi 
lokalisasi yang tinggi, tetapi rangkaian dua tahap membuat latensi inferensi meningkat 
sehingga sulit memenuhi kebutuhan \textit{real-time} pada perangkat berdaya rendah. Untuk 
sistem yang memproses aliran kamera secara terus-menerus, latensi yang tinggi menjadi 
hambatan utama.

Pendekatan berbasis \textit{transformer} yang diwakili RT-DETR menghilangkan tahap 
\textit{non-maximum suppression} (NMS) dan menjadikan proses deteksi bersifat 
\textit{end-to-end} \cite{zhao2024rtdetr}. Pada skala server, RT-DETR menawarkan 
keseimbangan kecepatan dan akurasi yang kuat. Varian RT-DETR-R50 mencapai 53,1\% AP pada 
COCO \textit{val2017} dengan 108 FPS pada GPU T4, mengungguli detektor YOLO sekelasnya 
dalam dua metrik tersebut \cite{zhao2024rtdetr}. Keterbatasannya muncul pada sisi 
efisiensi model. Varian terkecil RT-DETR, yaitu RT-DETR-R18, memiliki sekitar 20 juta 
parameter dan tidak menyediakan tier ringan setara kelas \textit{nano}. Selain itu, 
detektor berbasis \textit{transformer} cenderung membutuhkan data pelatihan yang lebih 
banyak untuk konvergen, sehingga kurang sesuai untuk \textit{dataset} kustom berukuran 
terbatas.

Pendekatan \textit{one-stage} berbasis CNN yang diwakili YOLO memprediksi kelas dan kotak 
pembatas dalam satu kali proses maju, sehingga inferensinya cepat \cite{redmon2016yolo}. 
Keluarga YOLO menyediakan rentang ukuran model yang lebar, mulai dari tier \textit{nano} 
dengan sekitar 2,6 juta parameter hingga varian terbesar 
\cite{terven2023comprehensive, khanam2024yolov11}. Pada sisi akurasi, varian terbesar 
YOLO11 mencapai 54,7\% mAP pada COCO dengan jumlah parameter yang lebih sedikit dibanding 
varian transformer setara, sehingga anggapan bahwa YOLO selalu kalah akurat terhadap 
\textit{transformer} tidak lagi berlaku mutlak. Ketersediaan tier ringan membuat YOLO 
dapat dijalankan pada perangkat terbatas, sementara ekosistem dan bobot 
\textit{pretrained} yang matang memudahkan pelatihan ulang pada \textit{dataset} kustom.

\begin{table}[h]
    \centering
    \caption{Perbandingan pendekatan deteksi objek terhadap kebutuhan sistem pemantauan.}
    \label{tab:perbandingan_metode}
    \begin{tabular}{|p{3.0cm}|p{3.0cm}|p{3.0cm}|p{3.0cm}|}
        \hline
        \textbf{Aspek} & \textbf{Faster R-CNN} & \textbf{RT-DETR} & \textbf{YOLO (v8/v11)} \\
        % \multicolumn{1}{|c|}{\textbf{Aspek}} & \multicolumn{1}{|c|}{\textbf{Faster R-CNN}} & \multicolumn{1}{|c|}{\textbf{RT-DETR}} & \multicolumn{1}{|c|}{\textbf{YOLO (v8/v11)}} \\   
        \hline \hline
        Pendekatan & \textit{Two-stage} CNN & \textit{Transformer} \textit{end-to-end} & \textit{One-stage} CNN \\
        \hline
        Kecepatan inferensi & Rendah & Tinggi & Tinggi \\
        \hline
        Varian teringan & Tidak ada tier ringan & $\approx$20 juta parameter & $\approx$2,6 juta parameter \\
        \hline
        Kesesuaian perangkat & Kurang (latensi tinggi) & GPU server & Perangkat terbatas \\
        terbatas & & & \\
        \hline
        Kebutuhan data latih & Sedang & Tinggi & Sedang \\
        \hline
    \end{tabular}
\end{table}

Mengacu pada empat tolok ukur permasalahan, YOLO memberikan kesesuaian terbaik. Faster 
R-CNN tersisih karena latensi yang tinggi, sedangkan RT-DETR, meski unggul pada skala 
server, tidak menyediakan tier ringan dan lebih menuntut data pelatihan. Penting dicatat 
bahwa keputusan ini bukan karena YOLO unggul secara mutlak di semua kondisi, melainkan 
karena ketersediaan tier ringan, efisiensi parameter, dan kematangan ekosistem YOLO 
selaras dengan kendala \textit{real-time}, perangkat terbatas, dan \textit{dataset} 
kustom yang menjadi ciri penelitian ini.

\subsection{Pemilihan Varian YOLO untuk Penelitian}
\label{subsec:pemilihan_varian}

Setelah keluarga YOLO ditetapkan, langkah berikutnya adalah memilih versi yang 
dibandingkan. Penelitian ini menggunakan YOLOv8 dan YOLO11 karena keduanya merupakan model 
unggulan resmi dari Ultralytics yang berbagi kerangka pelatihan, format data, alur 
augmentasi, dan antarmuka pemrograman yang identik 
\cite{terven2023comprehensive, khanam2024yolov11}. Kesamaan ini membuat perbedaan 
performa di antara keduanya dapat diatribusikan ke arsitektur, bukan ke perbedaan 
implementasi \textit{framework}, sehingga perbandingan menjadi terkontrol secara 
metodologis.

Versi YOLO lain sengaja tidak dimasukkan untuk menjaga kontrol tersebut. YOLOv9 
dikembangkan dari basis kode YOLOv5 dan diposisikan sebagai model eksperimental, 
sedangkan YOLOv10 dikembangkan oleh kelompok peneliti yang berbeda dengan mekanisme 
pelatihan tanpa NMS. Versi yang lebih baru seperti YOLOv12 belum memiliki kematangan 
tooling dan dukungan yang setara. Memasukkan versi-versi tersebut ke dalam perbandingan 
akan menambahkan variabel pengganggu berupa perbedaan asal basis kode dan filosofi desain. 
Membatasi perbandingan pada dua model unggulan Ultralytics yang setara justru menghasilkan 
perbandingan yang lebih valid.

Empat ukuran model, yaitu \textit{nano} (n), \textit{small} (s), \textit{medium} (m), dan 
\textit{large} (l), diuji pada tiap versi untuk memetakan keterkaitan antara ukuran model, 
akurasi, dan kecepatan yang relevan dengan kebutuhan \textit{real-time}. Varian terbesar 
(x) dikecualikan karena keterbatasan perangkat pelatihan dan deteksi serta tidak praktis 
untuk penerapan \textit{real-time}, dan pengecualian ini dinyatakan sebagai batasan 
penelitian.

Berdasarkan analisis di atas, metode yang dipilih adalah studi perbandingan YOLOv8 dan 
YOLO11 pada empat ukuran model, dilatih dan dievaluasi pada \textit{dataset} kustom 
kawasan Malioboro. Model dengan keseimbangan akurasi dan kecepatan terbaik akan dipilih 
untuk diterapkan pada sistem pemantauan. Rincian tahapan pelatihan, konfigurasi, dan 
evaluasi diuraikan pada Bab III.
